% Zone „Kennst du schon" – Nr. 1–12
\setcounter{aufgabe}{0}
\einheitenkopf[zone][Kennst du schon]{Das kennst du schon}

\begin{aufgabe}{Ich kann Dezimalzahlen multiplizieren und das Ergebnis runden.}
Rechne.
\begin{teilezwei}
\tz $4 \cdot 1{,}2$ \leerfeld & \tz $30 \cdot 0{,}9$ \leerfeld \\
\tz $250 \cdot 1{,}04$ \leerfeld & \tz $0{,}8 \cdot 0{,}8$ \leerfeld \\
\tz $7{,}5 \cdot 1{,}12$ \leerfeld & \\
\anweisung{Runde auf eine Nachkommastelle.}
\tz $45 \cdot 1{,}15 \cdot 1{,}15$ \leerfeld & \tz $1{,}8 \cdot 2{,}35$ \leerfeld \\
\end{teilezwei}
\end{aufgabe}

\begin{aufgabe}{Ich kann Punkte in ein Koordinatensystem eintragen.}
Trage die Punkte ein und beschrifte sie.
\begin{teile}
\teil $A(1\,|\,2)$
\teil $B(4\,|\,5)$
\teil $C(2{,}5\,|\,3{,}5)$
\end{teile}

\begin{ksys}[xmin=0,xmax=6,ymin=0,ymax=6]
\end{ksys}

\end{aufgabe}

\begin{aufgabe}{Ich kann Punkte eintragen, wenn eine Achse in größeren Schritten eingeteilt ist.}
Trage die Punkte ein und beschrifte sie.
\begin{teile}
\teil $D(1\,|\,100)$
\teil $E(3\,|\,250)$
\teil $F(4\,|\,375)$
\end{teile}

\begin{ksys}[xmin=0,xmax=5,ymin=0,ymax=400,ystep=50]
\end{ksys}

\end{aufgabe}

\begin{aufgabe}{Ich kann eine passende Einteilung für eine Achse wählen.}
Die senkrechte Achse hat eine bestimmte Zahl von Kästchen. Kreuze an, welche Einteilung je Kästchen die Achse am besten ausnutzt.
\begin{teile}
\teil Werte bis 40, 8 Kästchen: \quad \kreuz{1}\kreuz{5}\kreuz{10}
\teil Werte bis 900, 10 Kästchen: \quad \kreuz{50}\kreuz{100}\kreuz{200}
\teil Werte bis 2000, 10 Kästchen: \quad \kreuz{100}\kreuz{200}\kreuz{500}
\teil Werte bis 3,5, 7 Kästchen: \quad \kreuz{0,5}\kreuz{1}\kreuz{2}
\end{teile}
\end{aufgabe}

\begin{aufgabe}{Ich kann einen Wert um einen Prozentsatz erhöhen oder senken.}
Berechne den neuen Wert.
\begin{teile}
\teil 300\,€ werden um 10\,\% erhöht. \leerfeld[€]
\teil 50\,kg werden um 20\,\% gesenkt. \leerfeld[kg]
\teil 64\,€ werden um 5\,\% erhöht. \leerfeld[€]
\teil 180\,€ werden um 2,5\,\% gesenkt. \leerfeld[€]
\teil 45\,€ werden auf 80\,\% gesenkt. \leerfeld[€]
\teil 200\,€ werden erst um 10\,\% und danach noch einmal um 10\,\% erhöht. \leerfeld[€]
\end{teile}
\end{aufgabe}

\begin{aufgabe}{Ich kann eine Zahlenfolge fortsetzen, wenn immer gleich viel dazukommt.}
Setze die Folge um zwei Zahlen fort.
\begin{teile}
\teil 5, \ 8, \ 11, \ \leerfeld, \ \leerfeld
\teil 40, \ 35, \ 30, \ \leerfeld, \ \leerfeld
\teil 1,2; \ 1,5; \ 1,8; \ \leerfeld; \ \leerfeld
\end{teile}
\anweisung{Gib die Werte an.}
\begin{teile}
\teil Ein Taxi kostet 4\,€ Grundpreis und 2\,€ je Kilometer. Preis für 0\,km, 1\,km und 5\,km: \\ \leerfeld[€] \leerfeld[€] \leerfeld[€]
\teil An welcher Stelle trifft die Gerade $y = 3x + 5$ die senkrechte Achse? \leerfeld
\end{teile}
\end{aufgabe}

\begin{aufgabe}{Ich kann zwei Zahlen durcheinander teilen und das Ergebnis als Dezimalzahl angeben.}
Berechne den Quotienten.
\begin{teilezwei}
\tz $12 : 4$ \leerfeld & \tz $50 : 25$ \leerfeld \\
\tz $45 : 30$ \leerfeld & \tz $54 : 60$ \leerfeld \\
\anweisung{Teile den zweiten Wert durch den ersten.}
\tz von 80 auf 100: \leerfeld & \tz von 250 auf 200: \leerfeld \\
\end{teilezwei}
\end{aufgabe}

\begin{aufgabe}{Ich kann Prozentwert und Prozentsatz berechnen.}
Berechne den Prozentwert.
\begin{teilezwei}
\tz 10\,\% von 70\,€ \leerfeld[€] & \tz 25\,\% von 80\,kg \leerfeld[kg] \\
\tz 3\,\% von 450\,€ \leerfeld[€] & \\
\anweisung{Gib an, wie viel Prozent das sind.}
\end{teilezwei}
\begin{teile}
\teil 18 von 120 Schülern \leerfeld[\%]
\teil Ein Preis steigt von 50\,€ auf 56\,€. Um wie viel Prozent ist er gestiegen? \leerfeld[\%]
\end{teile}
\end{aufgabe}

\begin{aufgabe}{Ich kann ein Guthaben mit Zinseszins berechnen.}
Die Zinsen werden jedes Jahr zum Guthaben addiert. Berechne das Guthaben.
\begin{teile}
\teil 1000\,€ zu 2\,\% Zinsen, nach einem Jahr: \leerfeld[€]
\teil dasselbe Guthaben nach zwei Jahren: \leerfeld[€]
\teil 500\,€ zu 4\,\% Zinsen, nach drei Jahren: \leerfeld[€]
\end{teile}
\end{aufgabe}

\begin{aufgabe}{Ich kann Potenzen mit dem Taschenrechner berechnen.}
Berechne.
\begin{teilezwei}
\tz $2^5$ \leerfeld & \tz $10^3$ \leerfeld \\
\tz $1{,}5^3$ \leerfeld & \tz $(-2)^3$ \leerfeld \\
\anweisung{Runde auf zwei Nachkommastellen.}
\tz $0{,}9^4$ \leerfeld & \tz $3 \cdot 1{,}2^4$ \leerfeld \\
\end{teilezwei}
\end{aufgabe}

\begin{aufgabe}{Ich finde den Fehler in einer Zinseszinsrechnung.}
Tom legt 800\,€ für drei Jahre zu 5\,\% Zinsen an. Die Zinsen werden jedes Jahr zum Guthaben addiert. Tom rechnet so:
\rechnung{3 \cdot 5\,\% &= 15\,\% \\ 800\,€ \cdot 1{,}15 &= 920\,€}
Finde den Fehler, benenne ihn und korrigiere die Rechnung.
\end{aufgabe}

\begin{aufgabe}{Ich kann das Guthaben nach mehreren Jahren Zinseszins berechnen.}
Berechne das Guthaben.
\begin{teile}
\teil 1500\,€ zu 4\,\% Zinsen, nach zwei Jahren: \leerfeld[€]
\end{teile}
\end{aufgabe}
